\section{The difficulty and its local deficit}
\label{sec:difficulty-setup}

We now construct the integer that will measure progress along the
terminal chain of Proposition~\ref{prop:terminal-chain}.  Its definition
subtracts the ordinary families of codimension-two valuations from a
weighted discrepancy count and adds divisor-cycle ranks on the boundary
normalizations, following the framework of \cite[Section~2]{AHK}.
The ceiling weight used here can leave a negative correction where
several normalization branches meet.  We identify that correction
before proving the geometric estimate needed to bound it.

Work over $\mathbb C$ and suppose that a small klt program is infinite.
Use its terminal chain, with persistent prime labels $S_h$ and
nonincreasing coefficients $b_h\in[0,1)$, as constructed in
Section~\ref{sec:terminal}.  Retain zero-coefficient labels.  Make the
truncation in Proposition~\ref{prop:codim-two}, so that the fixed
coefficients are constant and, at every crepant junction, the varying
labels have positive coefficients and map downstairs to dimension at
most one.  Coefficients change only before each finite string of
small terminal flips.

In the big-boundary application of \cite{HanLiuZhuang2025}, uniform
Cartier-index bounds make discrepancies discrete, so the finitely many
extracted coefficients eventually stabilize
\cite[Lemma~3.4 and the proof of Theorem~3.1]{HanLiuZhuang2025}.
Here they may continue to decrease.  We clear denominators only for the
coefficients that have stabilized and control the rounding error from
the remaining labels by the branch inequality.

\subsection{Weights and finite local corrections}

Choose a positive integer \(N\) clearing the denominators of the
original downstairs boundary coefficients and of the fixed
coefficients on the terminal chain.  No denominator condition is
imposed on the varying coefficients.  Define
\begin{equation}
\label{eq:difficulty-weights}
 w(u)=\left\lceil N\max\{u,0\}\right\rceil,
 \qquad
 w_v(Y,\Theta)=w\bigl(2-a(v;Y,\Theta)\bigr),
 \qquad
 d_h=\sum_{j\geq 1}w\bigl(1-j(1-b_h)\bigr).
\end{equation}
By Lemma~\ref{lem:echoes}, $d_h$ is the default total echo weight
at a surface on the smooth one-label locus of $S_h$.
Only exceptional valuations over \(Y\) will enter the difficulty.
The notation \(w_v\) will be used when the pair is understood.

\begin{lemma}
\label{difficulty:preparation}
Each \(d_h\) is a finite nonnegative integer and is nonincreasing
along the terminal chain.  After discarding finitely many junctions,
all \(d_h\) are constant, including throughout every intervening
finite bridge.
\end{lemma}

\begin{proof}
Since \(b_h<1\), a summand in \eqref{eq:difficulty-weights} can be
positive only if \(j<1/(1-b_h)\).  Thus the sum is finite.  Each
summand is nonincreasing when \(b_h\) decreases.  At all subsequent
stages the possible indices \(j\) are bounded by the bound at the
initial stage under consideration.  Consequently \(d_h\) is a
nonincreasing sequence of nonnegative integers and is eventually
constant.  There are only finitely many labels.  Coefficients change
only at the coefficient-decrease move of a bridge, and stay unchanged
during its small flips.  Constancy at consecutive junctions therefore
gives constancy throughout the bridge as well.
\end{proof}

We henceforth make the truncation in
Lemma~\ref{difficulty:preparation}.  The constants \(d_h\) may depend
on the chosen infinite tail; a bound uniform over different programs
is neither asserted nor needed.

For an irreducible surface \(V\subset Y\), let \(r_{V,h}\) be the
number of prime divisors of \(S_h^\nu\) mapping onto \(V\), with
\(r_{V,h}=0\) when \(V\not\subset S_h\).  Here and below this is a
count of algebraic prime divisors over \(\kappa(\eta_V)\), not of
their geometric sheets.  Recall that
\(\rho_*(S_h^\nu)=\rho_2(S_h^\nu)\) is the rank of the span of
divisor cycle classes in \(H_4(S_h^\nu,\mathbb Q)\).

\begin{definition}
\label{def:difficulty}
For a terminal pair \((Y,\Theta)\) on the prepared chain, define
\begin{align}
\label{eq:difficulty}
 \mathcal D(Y,\Theta)
 ={}&\sum_{\operatorname{codim}_Y c_Y(v)\geq3}w_v
 \\
 &+\sum_{\substack{V\subset Y\text{ irreducible}\\\dim V=2}}
 \left(
       \sum_{c_Y(v)=V}w_v-\sum_h r_{V,h}d_h
 \right)
 +\sum_h d_h\rho_*(S_h^\nu).
 \notag
\end{align}
The correction in parentheses is performed at each individual surface
before summing over surfaces.
\end{definition}

Lemma~\ref{lem:echoes} makes this an integer defined by finite sums
of nonzero terms.  Indeed, the positive-weight contributions in
codimension at least three are finite, as are those at any one
surface center.  Outside finitely many exceptional surface centers,
a surface with a nonzero contribution lies generically in the
log-smooth locus on a single positive-coefficient label.  Its
contributing valuations are exactly the echoes with orders
\((j,1)\), whose log discrepancies are
\(1+j(1-b_h)\).  Their total weight is exactly \(d_h\), and the
local correction is zero.  Zero-coefficient labels have
\(d_h=0\) and do not affect this cancellation.  Thus
\eqref{eq:difficulty} is not a formal subtraction of two divergent
valuation sums.  Its local brackets, and initially its entire value
at an intermediate stage, are allowed to be signed.

\subsection{The local ceiling error}

\begin{lemma}
\label{lem:ceiling-error}
Let \(V\subset Y\) be an irreducible surface on a terminal model in
the prepared chain, and set
\[
 r_V^{\mathrm{var}}=\sum_{h\text{ varying}}r_{V,h}.
\]
Then
\begin{equation}
\label{eq:local-ceiling-bound}
 \sum_{c_Y(v)=V}w_v-\sum_h r_{V,h}d_h
 \ \geq\ -\max\{r_V^{\mathrm{var}}-1,0\}.
\end{equation}
The assertion does not require the boundary to have simple normal
crossings at \(\eta_V\).
\end{lemma}

\begin{proof}
The variety \(Y\) is terminal and hence smooth in codimension two.
Its local ring at \(\eta_V\) is therefore regular of dimension two.
Write \(m_h\) for the multiplicity there of \(S_h\), taking
\(m_h=0\) if the label does not contain \(V\).  The ordinary
blowup of the closed point of this local surface defines a valuation
\(v_1\) with
\begin{equation}
\label{eq:surface-first-blowup}
 a(v_1;Y,\Theta)=2-\sum_hm_hb_h>1.
\end{equation}
In particular \(\sum_hm_hb_h<1\), and the weight supplied by this
valuation is \(w(\sum_hm_hb_h)\).

We first justify the branch multiplicity inequality
\begin{equation}
\label{eq:branch-multiplicity}
 r_{V,h}\leq m_h.
\end{equation}
For a label containing \(V\), let \(R\) be the one-dimensional
local domain defined by its prime equation in this regular local
surface.  Its finite normalization \(\overline R\) is semilocal.
For a general parameter \(t\),
\[
 m_h=\operatorname{length}_R(R/tR)
     =\operatorname{length}_R(\overline R/t\overline R)
     =\sum_{\mathfrak q}
       [\kappa(\mathfrak q):\kappa(\eta_V)]
       \operatorname{ord}_{\mathfrak q}(t),
\]
where the sum runs over the branches of the normalization above the
closed point.  The middle equality follows by applying multiplication
by \(t\) to the finite-length module \(\overline R/R\): its kernel
and cokernel have equal length.  Each summand in the last expression
is a positive integer.  This proves
\eqref{eq:branch-multiplicity}, with the residue degrees included.

For fixed labels \(Nb_h\) is integral.  Apply the elementary
inequality
\(\lceil x_1+\cdots+x_r\rceil\geq
\sum_{\ell=1}^r\lceil x_\ell\rceil-(r-1)\) for \(r\geq1\)
to the varying branches, repeating \(Nb_h\) once for each branch
of that label.  When there are no varying branches the same comparison
is an equality for the fixed part.  Effectiveness and
\eqref{eq:branch-multiplicity} give
\begin{equation}
\label{eq:first-default-error}
 w\left(\sum_hm_hb_h\right)
 \geq \sum_h r_{V,h}w(b_h)
       -\max\{r_V^{\mathrm{var}}-1,0\}.
\end{equation}
Only the varying branches can cause this rounding error.

It remains to pay the terms with \(j\geq2\) in the defaults.
A positive such term requires \(b_h>1/2\).  By
\eqref{eq:surface-first-blowup}, at most one label through \(V\)
has this property, and its multiplicity must be one.  Its prime
equation is therefore a regular parameter, which we denote by \(x\);
complete it to regular parameters \((x,t)\).  In particular this
label has precisely one normalization branch at \(V\).
For each \(j\geq2\) with positive default, the monomial valuation
with orders \(v_j(x)=j\), \(v_j(t)=1\) is a normalized divisorial
valuation centered at \(V\).  These valuations are pairwise distinct
and different from \(v_1\).  If \(b_h\) is the coefficient of
\(x=0\), effectiveness of the other boundary components gives
\[
 a(v_j;Y,\Theta)
 =1+j-jb_h-\sum_{\ell\ne h}b_\ell v_j(S_\ell)
 \leq1+j(1-b_h).
\]
Thus \(w_{v_j}\geq w(1-j(1-b_h))\).  These distinct valuations
pay all higher default terms, while \(v_1\) gives
\eqref{eq:first-default-error}.  Summing proves
\eqref{eq:local-ceiling-bound}.
\end{proof}


\subsection{The geometric bound still needed}

At a crepant junction $p:Y\to X_i$, let
\[
 Q=\bigcup_{h\text{ varying}}S_h
\]
with its reduced structure.  For a surface $V\subset Q$, the number
$r_V^{\mathrm{var}}$ counts the prime divisors above $V$ on the disjoint
normalizations of the components of $Q$, and is at least one.
Surfaces outside $Q$ have no negative allowance in
Lemma~\ref{lem:ceiling-error}.  Summing that lemma therefore gives
\begin{equation}\label{eq:difficulty-local-deficit}
 \mathcal D(Y,\Theta)\geq
 \sum_{\operatorname{codim}_Y c_Y(v)\geq3}w_v
 +\sum_h d_h\rho_*(S_h^\nu)
 -\sum_{V\subset Q}(r_V^{\mathrm{var}}-1).
\end{equation}
If $Q$ is empty, this already proves nonnegativity.  Otherwise every
component of $Q$ has positive coefficient, so its default satisfies
$d_h\geq w(b_h)\geq1$.  Every valuation centered on a curve in $Q$
with log discrepancy below two likewise contributes at least one to
the first sum.  None of these valuations has a surface center.

Thus nonnegativity will follow if the total branch excess in the last
sum is bounded by the unweighted divisor-cycle ranks of the components
of $Q$, together with the number of these curve-centered valuations.
The next two sections prove precisely that bound, in
Theorem~\ref{thm:branch}.  Its geometric hypotheses hold here because
$\dim p(Q)\leq1$.  The lower bound is needed at the crepant junctions;
Section~\ref{sec:difficulty} will show that neither the coefficient
changes nor the small flips between them increase~$\mathcal D$.
